\singlespacing  


This report studied optimization strategies for static memory planning in TinyML inference, 
where minimizing the peak memory footprint of activation tensors is critical for deployment 
on resource-constrained embedded devices.

The experimental results show that the choice of memory allocation model significantly 
influences the achievable memory efficiency. Offset-based allocation methods, such as the 
Hole-aware Greedy algorithm, consistently achieve lower peak memory consumption than 
buffer-based approaches because tensors can be packed more densely in a continuous 
memory space.

Meta-heuristic optimization further improves the allocation quality by exploring different 
tensor ordering strategies. Among the evaluated approaches, the Genetic Algorithm provides 
the most effective memory reduction across most network topologies, although it requires 
longer runtime compared with simple greedy methods. In contrast, the Ping-pong strategy 
achieves the fastest execution time and remains competitive for simple sequential networks.

Overall, the results indicate that static memory planning performance depends both on the 
allocation model and the optimization strategy. Combining offset-based allocation with 
meta-heuristic search provides an effective approach for reducing peak memory usage in 
TinyML deployment scenarios.